\section{Conclusion}

This paper presented the design, simulation, and evaluation of an
Autonomous Smart Cart system built on ROS~2 Jazzy and Gazebo Harmonic.
The system integrates a UWB-based localisation pipeline, a 2D Kalman
filter, a proportional follow-me controller, and a LiDAR obstacle
safety layer to enable a differential-drive cart to autonomously track
a user within a simulated indoor retail environment.

Simulation results demonstrated that the implemented system meets all
primary functional objectives. The Kalman-smoothed UWB localisation
achieves a position error of $<$0.13\,m across the 1--4\,m operating
range, well within the 0.30\,m accuracy target. The follow-me controller
maintains a following distance of $1.0 \pm 0.15$\,m during both
straight-line and turning scenarios. The obstacle safety layer correctly
enforces three-zone proximity speed scaling and triggers emergency stops
at the specified 0.30\,m threshold, with no false stops caused by the
followed person due to the FOLLOW mode obstacle filter. End-to-end
control latency of $\approx$5\,ms satisfies the 30\,ms non-functional
requirement.

The modular ROS~2 architecture --- separating localisation, follow-me
control, navigation state management, and obstacle safety into
independent nodes connected via publish-subscribe topics --- proved
effective for iterative testing and independent tuning of each
subsystem. The two-stage velocity pipeline in particular provides a
clean separation between navigation intent and safety enforcement.

Key limitations of the current simulation include the use of a
teleop-controlled person model (rather than autonomous human motion),
a purely proportional controller (which introduces overshoot during
rapid directional changes), and the absence of SLAM-based global
localisation. These form the basis for the future improvements
described in Section~VIII, with computer-vision-based human tracking
and autonomous return-to-base identified as the highest-priority
second-generation upgrades.

Overall, the Gazebo simulation validates the core system architecture
and demonstrates that UWB-guided follow-me with layered LiDAR obstacle
avoidance is a viable and safe approach for autonomous indoor cart
systems. The codebase, launch configuration, and simulation world are
publicly available on GitHub for reproduction and further development.
